\Needspace{16\baselineskip}
\Article{13}

The factors on the right side of equation $[5]$ can also be written as follows:

\[
\begin{aligned}
r-r^{-1}&=-(r^{n-1}-r^{-n+1})\\
r^3-r^{-3}&=-(r^{n-3}-r^{-n+3})\\
r^5-r^{-5}&=-(r^{n-5}-r^{-n+5}),\;\text{etc.}
\end{aligned}
\]

up to

\[
\begin{aligned}
r^{n-2}-r^{-n+2}=-(r^2-r^{-2})
\end{aligned}
\]

The equation then takes this form:

\SourceLeaf{35}

\[
\begin{aligned}
W&=(-1)^{\frac12(n-1)}(r^2-r^{-2})(r^4-r^{-4})(r^6-r^{-6})\\&\qquad\dotsfour(r^{n-1}-r^{-n+1})
\end{aligned}
\]

Multiplying this equation by $[5]$ in its original form gives

\[
\begin{aligned}
W^2&=(-1)^{\frac12(n-1)}(r-r^{-1})(r^2-r^{-2})(r^3-r^{-3})\\&\qquad\dotsfour(r^{n-1}-r^{-n+1})
\end{aligned}
\]

where $(-1)^{\frac12(n-1)}$ is either $=+1$ or $=-1$, according as $n$ is of the form $4\mu+1$ or of the form $4\mu+3$.\quad Hence

\[
\begin{aligned}
W^2&=\pm r^{\frac12n(n-1)}(1-r^{-2})(1-r^{-4})(1-r^{-6})\\&\qquad\dotsfour(1-r^{-2(n-1)})
\end{aligned}
\]

But it is easy to see that $r^{-2}$, $r^{-4}$, $r^{-6}\dotsfour r^{-2n+2}$ give all the roots of the equation $x^n-1=0$ except the root $x=1$. The following identity must therefore hold for an indeterminate:

\[
\begin{aligned}
&(x-r^{-2})(x-r^{-4})(x-r^{-6})\dotsfour(x-r^{-2n+2})\\&\qquad=x^{n-1}+x^{n-2}+x^{n-3}+\text{etc.}+x+1
\end{aligned}
\]

Thus, putting $x=1$, we obtain

\[
\begin{aligned}
(1-r^{-2})(1-r^{-4})(1-r^{-6})\dotsfour(1-r^{-2n+2})=n
\end{aligned}
\]

and since evidently $r^{\frac12n(n-1)}=1$, our equation becomes

\[
\begin{aligned}
W^2=\pm n\  .\  .\  .\  .\  .\  .\  .\  .\  .\  .\ [6]
\end{aligned}
\]

Thus, in the case where $n$ is of the form $4\mu+1$, we have

\[
\begin{aligned}
W=\pm\surd n,\quad\text{and hence}\quad T=\pm\surd n,\quad U=0
\end{aligned}
\]

In the other case, where $n$ is of the form $4\mu+3$, we instead have

\[
\begin{aligned}
W=\pm i\surd n,\quad\text{and so}\quad T=0,\quad U=\pm\surd n
\end{aligned}
\]

\Needspace{13\baselineskip}
\Article{14}

The method of the preceding article gives only the absolute value of the sums $T$, $U$, and leaves it uncertain whether $T$ in the first case and $U$ in the second should be set $=+\surd n$ or $=-\surd n$.\quad But, at least for the case where $k=1$, this can be decided from equation $[5]$ as follows.\quad Since, for $k=1$,

\SourceLeaf{36}

\[
\begin{aligned}
r-r^{-1}&=2i\sin\omega\\
r^3-r^{-3}&=2i\sin3\omega\\
r^5-r^{-5}&=2i\sin5\omega\;\text{etc.}
\end{aligned}
\]

that equation becomes

\[
\begin{aligned}
W=(2i)^{\frac12(n-1)}\sin\omega\sin3\omega\sin5\omega\ldots\sin(n-2)\omega
\end{aligned}
\]

Now, in the case where $n$ is of the form $4\mu+1$, the sequence of odd numbers

\[
\begin{aligned}
1,\;3,\;5,\;7\dotsfour\tfrac12(n-3),\quad\tfrac12(n+1)\dotsfour(n-2)
\end{aligned}
\]

contains $\frac14(n-1)$ numbers smaller than $\frac12n$, to which positive sines evidently correspond. The remaining $\frac14(n-1)$, on the other hand, are greater than $\frac12n$, and negative sines correspond to them. The product of all the sines must therefore equal a positive quantity multiplied by $(-1)^{\frac14(n-1)}$. Thus $W$ will equal a positive real quantity multiplied by $i^{n-1}$, or by $1$, since $i^4=1$ and $n-1$ is divisible by $4$: that is, $W$ will be real and positive. Hence we must have

\[
\begin{aligned}
W=+\surd n,\quad T=+\surd n
\end{aligned}
\]

In the other case, where $n$ is of the form $4\mu+3$, in the sequence of odd numbers

\[
\begin{aligned}
1,\;3,\;5,\;7\dotsfour\tfrac12(n-1),\quad\tfrac12(n+3)\dotsfour(n-2)
\end{aligned}
\]

the first $\frac14(n+1)$ will be smaller than $\frac12n$, but the remaining $\frac14(n-3)$ will be greater.\quad Thus, among the sines of the arcs $\omega$, $3\omega$, $5\omega\dotsfour(n-2)\omega$, there will be $\frac14(n-3)$ negative ones, and so $W$ will be the product of $i^{\frac12(n-1)}$, a positive real quantity and $(-1)^{\frac14(n-3)}$. The third factor is $=i^{\frac12(n-3)}$, which, multiplied by the first, gives $i^{n-2}=i$, since $i^{n-3}=1$.\quad We must therefore have

\[
\begin{aligned}
W=+i\surd n,\quad\text{and}\quad U=+\surd n
\end{aligned}
\]

\Needspace{15\baselineskip}
\Article{15}

We shall now show how the same conclusions can be derived from the series considered in Article $9$.\quad In equation $[4]$, replace $x^{\frac12}$ by $-y^{-1}$; we then obtain

\SourceLeaf{37}

\[
\begin{aligned}
&1-y^{-1}\frac{1-y^{-2m}}{1-y^{-2}}+y^{-2}\frac{(1-y^{-2m})(1-y^{-2m+2})}{(1-y^{-2})(1-y^{-4})}\\&\qquad-y^{-3}\frac{(1-y^{-2m})(1-y^{-2m+2})(1-y^{-2m+4})}{(1-y^{-2})(1-y^{-4})(1-y^{-6})}+\text{etc.}
\end{aligned}
\]

up to the term in position $m+1$

\[
\begin{aligned}
&=(1-y^{-1})(1+y^{-2})(1-y^{-3})(1+y^{-4})\\&\qquad\ldots(1\pm y^{-m})\  .\  .\  .\  .\ [7]
\end{aligned}
\]

If we take for $y$ a proper root of the equation $y^n-1=0$, say $r$, and at the same time put $m=n-1$, we obtain

\[
\begin{aligned}
\frac{1-y^{-2m}}{1-y^{-2}}&=\frac{1-r^2}{1-r^{-2}}=-r^2\\
\frac{1-y^{-2m+2}}{1-y^{-4}}&=\frac{1-r^4}{1-r^{-4}}=-r^4\\
\frac{1-y^{-2m+4}}{1-y^{-6}}&=\frac{1-r^6}{1-r^{-6}}=-r^6
\end{aligned}
\]

up to

\[
\begin{aligned}
\frac{1-y^{-2}}{1-y^{-2m}}=\frac{1-r^{2n-2}}{1-r^{-2n+2}}=-r^{2n-2}
\end{aligned}
\]

Here we should note that none of the denominators $1-r^{-2}$, $1-r^{-4}$, etc. becomes $=0$.\quad Equation $[7]$ therefore takes this form:

\[
\begin{aligned}
&1+r+r^4+r^9+\text{etc.}+r^{(n-1)^2}\\&\qquad=(1-r^{-1})(1+r^{-2})(1-r^{-3})\dotsfour(1+r^{-n+1})
\end{aligned}
\]

Multiplying, on the right side of this equation, the first term by the last, the second by the next to last, etc., we obtain

\[
\begin{aligned}
(1-r^{-1})(1+r^{-n+1})&=r-r^{-1}\\
(1+r^{-2})(1-r^{-n+2})&=r^{n-2}-r^{-n+2}\\
(1-r^{-3})(1+r^{-n+3})&=r^3-r^{-3}\\
(1+r^{-4})(1-r^{-n+4})&=r^{n-4}-r^{-n+4}\;\text{etc.}
\end{aligned}
\]

\Needspace{8\baselineskip}
It is easy to see that these partial products combine to form the product

\[
\begin{aligned}
&(r-r^{-1})(r^3-r^{-3})(r^5-r^{-5})\\&\qquad\dotsfour(r^{n-4}-r^{-n+4})(r^{n-2}-r^{-n+2})
\end{aligned}
\]

which will therefore be

\[
\begin{aligned}
=1+r+r^4+r^9+\text{etc.}+r^{(n-1)^2}=W
\end{aligned}
\]

This equation is identical with equation $[5]$, derived in Article $12$ from the first series; the remaining arguments are then established in the same way as in Articles $13$ and $14$.

\SourceLeaf{38}

\Needspace{16\baselineskip}
\Article{16}

We turn to the other case, where $n$ is an even number.\quad First let $n$ be of the form $4\mu+2$, or twice an odd number. It is clear that the numbers $\frac14nn$, $(\frac12n+1)^2-1$, $(\frac12n+2)^2-4$, etc., or in general $(\frac12n+\lambda)^2-\lambda\lambda$, give odd quotients when divided by $\frac12n$, and are therefore congruent, modulo $n$, to $\frac12n$.\quad It follows that, if $r$ is a proper root of the equation $x^n-1=0$, and hence $r^{\frac12n}=-1$, we have

\[
\begin{aligned}
r^{(\frac12n)^2}&=-1\\
r^{(\frac12n+1)^2}&=-r\\
r^{(\frac12n+2)^2}&=-r^4\\
r^{(\frac12n+2)^2}&=-r^9\;\text{etc.}
\end{aligned}
\]

Thus, in the series

\[
\begin{aligned}
1+r+r^4+r^9+\text{etc.}+r^{(n-1)^2}
\end{aligned}
\]

the term $r^{(\frac12n)^2}$ cancels the first, the next cancels the second, etc., so that

\[
\begin{aligned}
W=0,\quad T=0,\quad U=0
\end{aligned}
\]

\Needspace{16\baselineskip}
\Article{17}

The case remains where $n$ is of the form $4\mu$, or divisible by four.\quad Here, in general, $(\frac12n+\lambda)^2-\lambda\lambda$ is divisible by $n$, and therefore

\[
\begin{aligned}
r^{(\frac12n+\lambda)^2}=r^{\lambda\lambda}
\end{aligned}
\]

Thus, in the series

\[
\begin{aligned}
1+r+r^4+r^9+\text{etc.}+r^{(n-1)^2}
\end{aligned}
\]

the term $r^{(\frac12n)^2}$ equals the first, the next equals the second, etc., so that

\[
\begin{aligned}
W=2(1+r+r^4+r^9+\text{etc.}+r^{(\frac12n-1)^2})
\end{aligned}
\]

Now suppose that in equation $[7]$ of Article 15 we put $m=\frac12n-1$, and take for $y$ a proper root of the equation $y^n-1=0$, say $r$.\quad Then, just as in Article 15, the equation takes the following form:

\[
\begin{aligned}
&1+r+r^4+\text{etc.}+r^{(\frac12n-1)^2}\\&\qquad=(1-r^{-1})(1+r^{-2})(1-r^{-3})\dotsfour(1-r^{-\frac12n+1})
\end{aligned}
\]

\SourceLeaf{39}

or

\[
\begin{aligned}
W&=2(1-r^{-1})(1+r^{-2})(1-r^{-3})(1+r^{-4})\\&\qquad\ldots(1-r^{-\frac12n+1})\quad\dotstwo\quad[8]
\end{aligned}
\]

\Needspace{10\baselineskip}
Further, since $r^{\frac12n}=-1$, and therefore

\[
\begin{aligned}
1+r^{-2}&=-r^{\frac12n-2}(1-r^{-\frac12n+2})\\
1+r^{-4}&=-r^{\frac12n-4}(1-r^{-\frac12n+4})\\
1+r^{-6}&=-r^{\frac12n-6}(1-r^{-\frac12n+6})\;\text{etc.}
\end{aligned}
\]

and the product of the factors $-r^{\frac12n-2}$, $-r^{\frac12n-4}$, $-r^{\frac12n-6}$, etc., up to $-r^2$ is $=(-1)^{\frac14n-1}\,r^{\frac1{16}nn-\frac14n}$, the preceding equation can also be written as

\[
\begin{aligned}
W&=2(-1)^{\frac14n-1}\,r^{\frac1{16}nn-\frac14n}(1-r^{-1})(1-r^{-2})\\&\qquad(1-r^{-3})(1-r^{-4})\dotsfour(1-r^{-\frac12n+1})
\end{aligned}
\]

\Needspace{10\baselineskip}
Since we have

\[
\begin{aligned}
1-r^{-1}&=-r^{-1}(1-r^{-n+1})\\
1-r^{-2}&=-r^{-2}(1-r^{-n+2})\\
1-r^{-3}&=-r^{-3}(1-r^{-n+3})\;\text{etc.}
\end{aligned}
\]

we obtain

\[
\begin{aligned}
&(1-r^{-1})(1-r^{-2})(1-r^{-3})\ldots(1-r^{-\frac12n+1})\\
&\quad=(-1)^{\frac12n-1}\,r^{-\frac18nn+\frac14n}(1-r^{-\frac12n-1})\\&\qquad(1-r^{-\frac12n-2})(1-r^{-\frac12n-3})\ldots(1-r^{-n+1})
\end{aligned}
\]

and hence

\[
\begin{aligned}
W&=2(-1)^{\frac34n-2}\,r^{-\frac1{16}nn}(1-r^{-\frac12n-1})\\&\qquad(1-r^{-\frac12n-2})(1-r^{-\frac12n-3})\dotsfour(1-r^{-n+1})
\end{aligned}
\]

Multiplying this value of $W$ by the one found earlier, and introducing the factor $1-r^{-\frac12n}$ on both sides, gives

\[
\begin{aligned}
(1-r^{-\frac12n})W^2&=4(-1)^{n-3}r^{-\frac14n}(1-r^{-1}\mathord{.})\\&\qquad(1-r^{-2})(1-r^{-3})\dotsfour(1-r^{-n+1})
\end{aligned}
\]

But we have

\[
\begin{aligned}
1-r^{-\frac12n}&=2\\
(-1)^{n-3}&=-1\\
r^{-\frac14n}&=-r^{\frac14n}
\end{aligned}
\]

\[
\begin{aligned}
(1-r^{-1})(1-r^{-2})(1-r^{-3})\dotsfour(1-r^{-n+1})=n
\end{aligned}
\]

\Needspace{10\baselineskip}
From this we finally conclude

\SourceLeaf{40}

\[
\begin{aligned}
W^2=2r^{\frac14n}n\  .\  .\  .\  .\  .\  .\  .\  .\  .\ [9]
\end{aligned}
\]

Now it is easy to see that $r^{\frac14n}$ is either $=+i$ or $=-i$, according as $k$ is of the form $4\mu+1$ or of the form $4\mu+3$.\quad And since

\[
\begin{aligned}
2i=(1+i)^2,\quad -2i=(1-i)^2
\end{aligned}
\]

in the case where $k$ is of the form $4\mu+1$, we have

\[
\begin{aligned}
W=\pm(1+i)\surd n,\quad\text{and so}\quad T=U=\pm\surd n
\end{aligned}
\]

but in the other case, where $k$ is of the form $4\mu+3$,

\[
\begin{aligned}
W=\pm(1-i)\surd n,\quad\text{and so}\quad T=-U=\pm\surd n
\end{aligned}
\]
